\section{Projektas}

\subsection{Reikalavimų specifikacija}

\subsubsection{Komercinė specifikacija}

\subsubsection{Sistemos funkcijos}

\subsubsection{Vartotojo sąsajos specifikacija}

\subsubsection{Realizacijai keliami reikalavimai}

\subsubsection{Techninė specifikacija}

\subsection{Projektavimo metodai}

\subsubsection{Projektavimo valdymas ir eiga}

\subsubsection{Projektavimo technologija}

\subsubsection{Programavimo kalbos, derinimo, automatizavimo priemonės, operacinės sistemos}

\subsection{Sistemos projektas}

\subsubsection{Statinis sistemos vaizdas}

% TODO (question)
% Paklausti ar su italic pavadinimus ar su kabutėmis.
Pasirinkta dirbti su dvejom skirtingom \glsdisp{simulator}{imituoklėm}: atviros prieigos \textit{Veripool Verilator} ir komercine \gls{eda} „Cadence Xcelium“. Dėl esminių įrankių skirtumų privaloma palaikyti dvi atskiras sistemas, kurios leidžia dirbti su viena iš anksčiau išvardytų \glsdisp{simulator}{imituoklių} ir sukonfigūruoja projekto naudojamus įrankius. „Veripool Verilator“ \glsdisp{simulator}{imituoklės} darbo aplinka yra universali, nes kartu su ja yra naudojami kiti atviros prieigos įrankiai \pav{\ref{fig:uml_deplayement_verilator}}. „Cadence Xcelium“ \glsdisp{simulator}{imituoklės} darbo aplinka yra pririšta prie „Rocky Linux“ operacinės sistemos \pav{\ref{fig:uml_deplayement_xcelium}}, nes šie \gls{eda} įrankiai yra komerciniai ir jų veikimui yra reikalinga licenzija bei užtikrina pilną palaikymą visiems \gls{eda} įrankiams tik ant „Red Hat Enterprise Linux“ operacinės sistemos \cite{cadence_supported_platforms} ir ja pagrįstų skirstinių, tokie kaip „Rocky Linux“ \cite{cadence_support_roadmap}. Būtent dėl šių priežasčių karkasas turi dvi atskiras darbo aplinkas, kurių valdymas yra per šakninį „GNU Make“ surinkimo sistemos (angl. \textit{build system}) skripto failą.

\begin{figure}[H]
	\includesvg[width=11.22cm]{images/uml_deployement_verilator.svg}
	\caption{Darbo aplinkos su atviros prieigos įrankiu „Veripool Verilator“ diegimo diagrama}
	\label{fig:uml_deplayement_verilator}
\end{figure}

\begin{table}[H]
	\caption{\textit{Verilator} darbo aplinkos naudojami įrankiai}
	\label{tab:docker_dependencies_verilator}
	\begin{tabularx}{\linewidth}{|>{\raggedright\arraybackslash}p{4cm}|X|}
		\hline
		\thead[l]{Įrankiai}                       & \thead[l]{Paskirtis}                                                                                                                        \\
		\hline
		\textit{Verilator v5.044}                 & Atvirosios prieigos aparatinės įrangos aprašymo kalbos \gls{hdl} \glsdisp{simulator}{imituoklė}                                                                      \\
		\hline
		\textit{GTKWave Analyzer}                 & Skaitmeninių signalų laikinių diagramų \angl{waveforms} peržiūros įrankis.                                                                  \\
		\hline
		\textit{GNU Build Essentials}             & C++ kompiliatorius ir surinkimo įrankiai, reikalingi \textit{Verilator} modelių generavimui iš Verilog/SystemVerilog \gls{rtl} \gls{mcu} pradinio kodo. \\
		\hline
		\textit{Pagalbinės Verilator bibliotekos} & \textit{Verilator} veikimui ir leksinei analizei reikalingos bibliotekos.                                                                   \\
		\hline
		\textit{\gls{gui}}                              & Vartotojo grafinei sąsajai naudojamos paslaugų programos \angl{utilities}, jog būtų galima atvaizduoti \textit{GTKWave} langus.             \\
		\hline
	\end{tabularx}
\end{table}

Aukščiau pateiktoje \gls{uml} diegimo diagramoje \pav{\ref{fig:uml_deplayement_verilator}} yra detaliai atvaizduota „Veripool Verilator“ darbo aplinkos sistema ir jos diegimas. Naudojamas \textit{Docker} konteinerizavimo variklis, kuris leidžia beveik visiškai atskirti projekto darbo aplinką nuo pagrindinės \angl{host} operacinės sistemos. Tai reiškia, jog jei yra palaikomas \textit{Docker} variklis ant operacinės sistemos platformos, tokios kaip: \textit{Linux}, \textit{MacOS}, \textit{Windows} \cite{docker_docs}, visa projekto darbo aplinka ir įrankiai prieinami yra prienami \glsdisp{container}{konteineryje}, nepriklausant nuo platformos operacinės sistemos ir jos specifinių konfigūracijų, nustatymų. Antras įrankis, kuris privalo būti įdiegtas pagrindinėje sistemoje yra \textit{GNU Make} surinkimo sistema, kuri valdo darbo aplinkos diegimo eigą. \textit{GNU Make} taip pat yra prieinamas įvairioms platformoms: \textit{Linux}, \textit{MacOS}, \textit{Windows}, \textit{BSD} \cite{gnu_make_supported_platforms}. \Glsdisp{container} yra suglaudinta \angl{compressed} \textit{Ubuntu} operacinė sistema, kurioje yra įdiegiami visi reikalingi įrankiai naudojant \textit{APT} paketų tvarkytuvę \angl{package manager}. \Glsdisp{container}{Konteinerio} aprašą naudoja automatizuotas \textit{GNU Make} skriptas, kuris įdiegia ir sutvarko \textit{Verilator} \glsdisp{simulator}{imituoklei} reikalingas priklausomybes \lent{\ref{tab:docker_dependencies_verilator}} ir projekto darbo eigai reikalingus įrankius ir priklausomybės \lent{\ref{tab:docker_dependencies}}.

\begin{figure}[H]
	\includesvg[width=11.22cm]{images/uml_deployement_xcelium.svg}
	\caption{Darbo aplinkos su komerciniu \gls{eda} įrankiu „Cadence Xcelium“ diegimo diagrama}
	\label{fig:uml_deplayement_xcelium}
\end{figure}

\begin{table}[H]
	\caption{\textit{Xcelium} darbo aplinkos specifiniai paketai ir įrankiai}
	\label{tab:docker_dependencies_xcelium}
	\begin{tabularx}{\linewidth}{|>{\raggedright\arraybackslash}p{4.5cm}|X|}
		\hline
		\thead[l]{Paketai / Įrankiai}                     & \thead[l]{Paskirtis}                                                                                                                                                                                   \\
		\hline
		% TODO (source)
		\textit{GCC kompiliatorius ir surinkimo įrankiai} & Būtini „Xcelium“ \glsdisp{simulator}{imituoklei} atliekant \gls{rtl} transliaciją, DPI-C (angl. \textit{Direct Programming Interface}) kodo kompiliavimą ir simuliacijos momentinių kopijų (angl. \textit{snapshots}) generavimą. \\
		\hline
		\textit{X11 \gls{gui} bibliotekos}                      & Reikalingos „SimVision“ laikinių diagramų peržiūros programos, atvaizdavimui.                                                                                                                          \\
		\hline
		% TODO (source)
		\textit{Suderinamumo bibliotekos}                 & C bibliotekos, užtikrinančios 32-bitų architektūros palaikymą 64-bitų sistemoje. Būtinos komercinių \gls{eda} įrankių licencijų valdymo posistemėms..                                                        \\
		\hline
	\end{tabularx}
\end{table}

Pateiktoje \gls{uml} diagramoje \pav{\ref{fig:uml_deplayement_xcelium}} yra pavaizduota komercinės \textit{Cadence Xcelium} \glsdisp{simulator}{imituoklės} darbo aplinkos sistema ir jos diegimas. Sistemos architektūra nemažai skiriasi nuo anksčiau aptartos \textit{Veripool Verilator} aplinkos \pav{\ref{fig:uml_deplayement_verilator}}. To priežastis yra ta, jog \textit{Cadence Xcelium} yra komercinė licencijuota \gls{eda} įranga ir licenzijos yra pririštos prie pačių įrankių vykdomųjų failų. Dėl šitos priežasties, \glsdisp{container}{konteineris} įrengia \angl{mount} \textit{Cadence} \gls{eda} įrankių aplankalą, esantį pačiame serveryje, su tik skaitymo \angl{read-only} prieiga. Šiuo būdu, konteinerizuota \textit{Python 3} testavimo aplinka gali kreiptis per \gls{vpi} sąsają tiesiai į pačio serverio \textit{Cadence Xcelium} vykdomąjį failą ir taip paleisti \glsdisp{simulator}{imituoklės} procesą \glsdisp{container}{konteinerio} viduje. Norint sėkmingai paleisti \glsdisp{simulator}{imituoklę} \glsdisp{container}{konteinerizuotoje aplinkoje}, reikia pateikti serveryje esančias licenzijas \glsdisp{container}{konteineriui } per \texttt{CDS\_LIC\_FILE} serverio \glsdisp{environment-variable}{aplinkos kintamąjį}. Vietoje \textit{Docker} variklio, \textit{Rocky Linux} turi savo alternatyvą \textit{Podman}, kuri yra įtraukta į sistemos įrankius \cite{rocky_linux_podman_guide}, šis variklis yra naudojamas \glsdisp{container}{konteinerių} kūrimui ir valdymui. Konteinerio aprašą taipogi naudoja tas pats automatizuotas \textit{GNU Make} skriptas, kuris įdiegia reikalingas priklausomybės \textit{Cadence Xcelium} ir \textit{SimVision} \gls{eda} įrankiams \lent{\ref{tab:docker_dependencies_xcelium}} bei perduoda licenzijas ir nurodo \gls{eda} įrankių fizinę vietą pačiame serveryje, taip pat sutvarko ir suderina pačio projekto eigai reikalingus įrankius ir priklausomybes \lent{\ref{tab:docker_dependencies}}.

\begin{table}[H]
	\caption{Bendri \textit{Verilator} ir \textit{Xcelium} konteinerių naudojami įrankiai}
	\label{tab:docker_dependencies}
	\begin{tabularx}{\linewidth}{|>{\raggedright\arraybackslash}p{4cm}|X|}
		\hline
		\thead[l]{Įrankiai}           & \thead[l]{Paskirtis}                                                                                        \\
		\hline
		\textit{GNU \gls{risc}-V Toolchain} & Kompiliatorius, skirtas ruošti C \glsdisp{program-build}{programinio kodo darinius} \gls{risc}-V 32 ir 64 bitų architektūroms.             \\
		\hline
		\textit{Python v3.12}         & \textit{Python} skriptų vykdymas, \gls{uvm} testavimo aplinkos valdymas.                                    \\
		\hline
		\textit{Sisteminiai įrankiai} & Kodo eilučių skaičiavimo, failų struktūros atvaizdavimo, kodo pataisos ir kiti projekto naudojami įrankiai. \\
		\hline
	\end{tabularx}
\end{table}

Abiejose konteinerizuotose darbo aplinkose kartu yra įdiegiami ir sukonfigūruojami bendri projekto darbo įrankiai \lent{\ref{tab:docker_dependencies}}. Konteineriai turi tiesioginę prieigą prie projekto failų. Kompiliatorius, skirtas C pirminiam programiniam tekstui \angl{source code}, naudoja pagrindinėje sistemoje esančius projekto poaplankio, skirto \gls{mcu} C kodui. Kompiliatoriaus darinių failai yra prieinami ir iš pačios sistemos ir iš konteinerizuotų darbo aplinkų. Kartu yra suderinami ir įdiegiami įrankiai, kurie naudojami pačio projekto, kaip \glsdisp{patch}{pataisos}, analizės, formatavimo ir kitos naudojamos programinės įrangos.

\begin{figure}[H]
	% TODO (refac)
	\includesvg[width=11.22cm]{images/uml_package.svg}
	\caption{Verifikavimo \glsdisp{framework}{karkaso} \gls{uml} paketu diagrama}
	\label{fig:uml_package}
\end{figure}

Karkasas yra pritaikytas naudoti su jau parašyta sistema kaip verifikacijos \glsdisp{plugin}{įskiepis}, kuris gali būti naudojamas keliuose \gls{mcu} projektuose, kurie gali turėti skirtingas implementacijas, svarbu tik tai, jog pagrindiniai protokolai ir funkcionalumai būtų palaikomi. Tai yra padaryta specialiai, jog būtų galima šį \glsdisp{framework}{karkasą} naudoti tarp skirtingų projekto stadijų, kur \gls{mcu} yra toliau modifikuojamas ir tikrinama, ar neregresavo anksčiau realizuoti funkcionalumai su naujais pakeitimais ar redagavimais. Pačio projekto struktūra detalizuoja \gls{uml} paketų diagrama \pav{\ref{fig:uml_package}}.

Projektas sudarytas iš šių esminių paketų:

\begin{itemize}
	\item Programinės įrangos paketas \texttt{sw} yra projekto programinė įranga, joje yra aprašyti pagrindiniai \gls{mcu} programiniai funkcionalumai C programavimo kalba. Šis paketas savyje turi pirminius failus, kurie yra naudojami \gls{mcu} programinės įrangos kompiliavimui ir darinių suformavimui, kurie yra įkeliami į \glsdisp{microprocessor}{mikrovaldiklio} instrukcijų (\textit{IMEM}) ir duomenų (\textit{DMEM}) atmintis imitavimo pradžios metu. Taip pat turi ir savo viduje esančius paketus:

	      \begin{itemize}
		      \item Bendrasis paketas \texttt{comm} {common} vykdo bendrų programinių priemonių paskirtį. Laiko failus, kurie yra bendri tarp visų \texttt{sw} vidinių paketų.
		      \item \gls{soc} paketas \texttt{soc} turi C antraštinius \angl{header} programinius failus, kuriuose yra bazinės struktūros \angl{structs}, kurių pagalba yra valdomi \glsdisp{microprocessor}{mikrovaldiklio} periferiniai įrenginiai. Apjungia programinę įrangą su aparatine įranga.
		      \item \gls{hal} paketas \texttt{hal} turi C programinius aprašus, kurie suteikia prieigą funkciniam periferinių įrenginių manipuliavimui, kaip duomenų rašymas, skaitymas.
		      \item Bibliotekos paketas \texttt{lib} \angl{library} suteikia įvairiausius papildomus sisteminius funkcionalumus, kaip eilučių \angl{string} struktūras ir operacijas su jomis.
		      \item Sistemos paketas \texttt{sys} \angl{system} susideda iš programinių C failų, kurie apibrėžia pačios sistemos veikimą ir ypatumus, kaip visuotinius kintamuosius \angl{global variables}, sistemines klaidas, pertraukčių \angl{iterrupts} apdorojimo apibrėžtis ir pan.
		      \item Įvesties ir išvesties paketas \texttt{io} \angl{Input-Output} aprašo \gls{cli} funkcionalumą ir suteikia įvesčių bei išvesčių funkcionalumus.
		      \item  Testų paketas \texttt{test} aprašo individualius testinius C programinius scenarijus, kurie yra skirti \gls{uvm} testavimo aplinkai ir yra keliami į \glsdisp{simulator}{imituoklės} imituojamo \glsdisp{microprocessor}{mikrovaldiklio} instrukcijų ir duomenų atmintis.
	      \end{itemize}

	\item Testavimo stendo \texttt{tb} paketas realizuoja projekto \gls{uvm} verifikavimo klases ir charakteristikas. Naudojami \gls{uvm} komponentai ir objektai aprašyti pagrindinius verifikacijos klases, kurios skirtos naudoti testų rašymui Python programavimo kalba. Viduje esantys paketai yra:

	      \begin{itemize}
		      \item Etaloninio modelio \texttt{ref} \angl{reference model} paketas, kuris realizuoja \glsdisp{microprocessor}{mikrovaldiklio} etaloninį modelį ir naudojamas gauti numatytus \angl{expected} rezultatus iš etaloninio modelio.
		      \item Aplinkos \texttt{env} \angl{environment} paketas, kuriame yra \glsdisp{microprocessor}{mikrovaldiklio} testavimo aplinka, skirta naudoti ir konfigūruoti \gls{uvm} testuose.
		      \item \gls{uvc} \texttt{uvc} paketas, jame yra sudėti \glsdisp{microprocessor}{mikrovaldiklio} \glsdisp{peripheral}{periferinių įrenginių} verifikavimo komponentai.
		      \item Virtualios sąsajos \texttt{vif} \angl{virtual interface} paketas suteikia sąsają su testuojamu \glsdisp{microprocessor}{mikrovaldiklio} \gls{rtl} viršutiniu \angl{top} \glsdisp{module}{moduliu}, naudojamas verifikacijos komponentuose ir testuose.
		      \item Konfigūracijos \texttt{cfg} \angl{configuration} paketas skirtas testavimo aplinkos specifiniam konfigūravimui tarp testų.
		      \item Sekų \texttt{seq} \angl{sequence} paketas įgyvendina nepriklausomas sekas, kurių pagalba vykdomas mikrovaldiklio verifikavimas testuose.
		      \item Testų \texttt{test} paketas kuriame yra keliami \gls{uvm} testai.
	      \end{itemize}

	\item Bendrosios aplinkos \texttt{env} paketas skirtas laikyti testuojamosios aplinkos specifikas, toks kaip programos dydis, pavadinimas, parametrai ir pan. Nustatomas simuliacijos \texttt{sim} paketo ir yra naudojamas \gls{uvm} testavimo stendo \texttt{tb} paketo.

	\item \gls{rtl} \texttt{rtl} paketas, kuriame yra aprašytas testuojamojo mikrovaldiklio \glsdisp{module}{modulis}, jis susideda iš jau aprašytų \glsdisp{module}{modulių}, esančių \texttt{mlab\_mcu} šakninio projekto \texttt{rtl} paketo.

	\item Skriptų \texttt{script} paketas kuriame yra automatizuotos eigos Bash kalba aprašyti pernaudojami \angl{reusable} scenarijai.

	\item Šakninio projekto \texttt{mlab\_mcu} \texttt{misc} paketas, kuriame yra duomenys ir aprašai, naudojami imitacijos vykdymo metu.

	\item Pataisos \texttt{patch} paketas suteikia pataisas šakniniam projektui, reikalingos tinkamam suderinamumui su verifikavimo \glsdisp{framework}{karkasu}.

	\item Konfigūracijos \texttt{conf} paketas, skirtas įvairiems konfigūraciniams failams, skirtų įvairių projekto aplinkų konfigūracijų taikymui.

	\item Imitacijos \texttt{sim} \angl{simulation} paketas, skirtas testų inicializavimui ir paleidimui \textit{Verilator} arba \textit{Xcelium} \glsdisp{simulator}{imituoklėse}.
\end{itemize}

\begin{figure}[H]
	\includesvg[width=11.22cm]{images/uml_package_tb.svg}
	\caption{\gls{uvm} testų \texttt{tb} su vidinėmis klasėmis paketų diagrama}
	\label{fig:uml_package_tb}
\end{figure}

\gls{uvm} testavimo paketas \texttt{tb} \pav{\ref{fig:uml_package_tb}}. Detalesnis jo vidinių paketų paskirčių aprašymas:

\begin{itemize}
	\item \glsdisp{golden-reference}{Etaloninio modelio} \texttt{ref} paketas -- \gls{uvm} komponentai ir objektai, kurie yra skirti tikrinti \gls{dut} elgseną transakcijų lygyje, juos naudoja \texttt{env} paketo \texttt{McuScoreboar} \gls{uvm} komponentas. 
	\item Aplinkos \texttt{env} paketas -- \gls{uvm} aplinkos komponentai, \texttt{McuEnv} yra konfigūruojamas globalia konfigūracija, komponentas pasirūpina visų vidinių komponentų sujungimu ir konfigūravimu ir yra pernaudojamas komponentas tarp visų testų.
	\item Testų \texttt{test} paketas -- visi aprašyti \gls{dut} testai, kuriuose naudojamos nepriklausomos sekos, visi testai paveldi iš bazinio \texttt{McuTest}.
	\item \gls{uvc} komponentų \texttt{uvc} paketas -- \gls{uart} ir \gls{gpio} \glsdisp{peripheral}{periferinių įrenginių}: įtaiso tvarkyklės \texttt{Driver}, monitoriai \texttt{Monitor}, sekų valdytojai \texttt{Sequencer} ir sąsajos \texttt{Interface}. Konfigūruojami agentai \texttt{Agent} yra naudojami kaip tarpininkai tarp \gls{dut} \glsdisp{peripheral}{periferijų} ir testinės \gls{mcu} aplinkos \texttt{McuEnv}.
	\item Virtualios \gls{mcu} sąsajos \texttt{vif} paketas -- tarpininkas tarp imituojamo \gls{dut} ir \textit{PyUVM}, per šią sąsają agentų komponentai komunikuoja su pačiu \gls{dut}.
	\item Nepriklausomų sekų ir jų elementų \texttt{seq} paketas -- tarp testų naudojamos sekos, kurie gali būti atsitiktiniai arba ranka aprašyti scenarijai, kurie perduoda transakcijas iš sekų \texttt{Sequence} per virtualų sekų valdytoją \texttt{McuVirtualSequencer}, kuris nukreipia į atitinkamą \glsdisp{peripheral}{periferijos} \gls{uvc} sekos valdytoją.
	\item Klaidų \texttt{errors} paketas -- su \glsdisp{verification}{verifikavimo} susijusių klaidų hierarchiją, kuri leidžia lengviau valdyti ir atrinkti įvairias testavimo metu iškilusias klaidas.
\end{itemize}

Bendram vaizdui, žemiau yra pateikta pilna \glsdisp{verification}{verifikacijos} \gls{uvm} objektų ir komponentų, \textit{PyVSC} atsitiktinių klasių ir pačių \textit{Python} \gls{uml} klasių diagrama \pav{\ref{fig:uml_class_tb_full}}.

\begin{figure}[H]
	\includesvg[width=11.22cm]{images/uvm_class_tb_full.svg}
	\caption{\gls{uvm} testų paketo \texttt{tb} pilna \gls{uml} klasių diagrama}
	\label{fig:uml_class_tb_full}
\end{figure}

% TODO (feat)
% class diagram of the UVM side (imagine, writing C++ instead of C).



\subsubsection{Dinaminis sistemos vaizdas}
